\section{Scriptural Date and Location}
\label{sec:date-location}
{\small Each passage the Mass reads or sings from Scripture, once, in
canonical order; beneath each, its historical judgment and critical horizon.\par}

\begin{dossiertable}
Offertory & Ex 24:4--5, compiled & Sinai: an altar at the foot of the mount,
with twelve titles for the twelve tribes of Israel (24:4) &
\chronodate{offertory}{\chronologyannotation{offertory}}\\
\cmidrule(lr){1-4}
\dossierevent{Narrated event: the ratifying of the covenant at Sinai, after
Moses has told the people all the words of the Lord and they have answered with
one voice (24:3); the altar is built ``rising in the morning'' (24:4), a
relative date within the narrative.}
\cmidrule(lr){1-4}
\dossierprose{The chronology record dates the narrated event only, and only
relatively; it holds no composition date for the Book of Exodus at these
verses, and none is given here. The antiphon is the liturgy's own compilation:
beside the altar, holocausts and victims of 24:4--5 it sets the verb of Moses'
consecration of the altar (Lev 8:11; Num 7:1), the evening sacrifice and odour
of sweetness of the daily offering (Ex 29:41; Ps 140:2), and the closing
\latin{in conspectu filiorum Israel} of 24:17, where the glory of the Lord is
seen on the mountain.}
\midrule
Communion & Ps 95:8--9 (Heb.\ 96) & The house of God and its courts; by the
title, the house built after the captivity, at Jerusalem &
\chronodate{communion}{\chronologyannotation{communion}}\\
\cmidrule(lr){1-4}
\dossierprose{Titled in the Vulgate \latin{Canticum ipsi David, quando domus
aedificabatur post captivitatem}, ``A canticle for David himself, when the house
was built after the captivity'' (v.~1). The title names David and a setting
after the exile together; the chronology record reads that setting against the
raising of the Temple from its ruins under Zorobabel after the return (the
\work{Catholic Encyclopedia}, ``Temple of Jerusalem''), and holds no Davidic era
for the psalm. For its composition the record gives only the modern critical
boundary for the whole Psalter, within which no single psalm can be dated
securely. The third date is not a date of the psalm: the record reads Ps~95,
with Pss~67, 96 and 97, as looking prophetically to Christ's coming, and the two
years it gives are two \work{Catholic Encyclopedia} articles' dates for the
Nativity.}
\midrule
Alleluia & Ps 101:16 (Heb.\ 102) & Sion, to be rebuilt and seen in glory
(vv.~14, 17) & \chronodate{alleluia}{\chronologyannotation{alleluia}}\\
\cmidrule(lr){1-4}
\dossierprose{Titled ``The prayer of the poor man, when he was anxious, and
poured out his supplication before the Lord'' (v.~1); the title names no author,
and the chronology record holds no traditional date for the psalm. Its lament
turns at verse 14 to Sion, on which the Lord will have mercy, and verse 17 gives
the reason for the nations' fear: ``the Lord hath built up Sion''. The date is
the Psalter's critical boundary only.}
\midrule
Gradual & Ps 121:1, 7 (Heb.\ 122); verse 1 is also the Introit's verse &
Jerusalem: the house of the Lord, and the city's courts, walls and towers
(vv.~1--7) & \chronodate{gradual}{\chronologyannotation{gradual}}\\
\cmidrule(lr){1-4}
\dossierprose{Titled in the Vulgate \latin{Canticum graduum}, one of the Songs
of Ascents (from Ps~119 to Ps~133 in the Missal's numbering), sung by pilgrims going up
to the house of the Lord; the Vulgate title names no author, and the chronology
record holds no traditional date for the psalm. Among its ancient interpreters
St John Chrysostom and Theodoret set it at the return from the Babylonian exile,
of a city whose walls and towers lay in ruins. The date is the Psalter's
critical boundary only.}
\midrule
Introit & Ecclus 36:18 (the antiphon, adapted); its verse, Ps 121:1, stands
with the Gradual & ``Jesus the son of Sirach, of Jerusalem'' (50:29), praying
for Jerusalem and Sion (36:14--16) &
\chronodate{introit}{\chronologyannotation{introit}}\\
\cmidrule(lr){1-4}
\dossierprose{Traditional attribution: Jesus the son of Sirach, whom the Greek
prologue names and whom the Greek and Latin text calls a man of Jerusalem (50:29); his
grandson rendered the book into Greek after coming into Egypt (\work{Catholic
Encyclopedia}, ``Ecclesiasticus'', which calls the authorship universally
assigned to him). The chapter is a prayer for Israel's deliverance, for
Jerusalem, and for the fulfilment of the prophets' words. The chronology record
dates the element's composition passage by passage. Ecclus~36:18: B.C.~190--170
or c.~B.C.~280, the composition of the Book of Ecclesiasticus under the
traditional Catholic profile, the interval preferred and the approximate year
its alternate, both from the \work{Catholic Encyclopedia}'s article
``Ecclesiasticus'' (1909). Ps~121:1: before c.~165~B.C., the Psalter's critical
composition boundary, already given with the Gradual. Neither date holds at
both loci, so the Date cell gives the element none. The antiphon's wording is
liturgical: at its opening and its close it is neither the Vulgate's nor the
Greek's.}
\midrule
% The two tiers of a dossier are read together; the sheet turns between
% dossiers, before the Gospel, rather than inside one.
\newpage
Gospel & Mt 9:1--8 & Matthew in Judea, for the Jews, before his departure from
Jerusalem. Event: ``his own city'' across the lake (9:1),
Capernaum for Chrysostom, Nazareth for Jerome &
\chronodate{gospel}{\chronologyannotation{gospel}}\\
\cmidrule(lr){1-4}
\dossierevent{Narrated event: the healing of the paralytic, after the crossing
from the country of the Gerasens (8:28--9:1) and before the call of Matthew
(9:9); the chronology record gives the episode no date.}
\cmidrule(lr){1-4}
\dossierprose{Traditional attribution: St Matthew the Apostle, whom the
Missal's heading names (\latin{secundum Matthaeum}). The Date column keeps the
\work{Catholic Encyclopedia}'s ranges as disputed alternatives, not equally
established estimates. The article ``Gospel of St.\ Matthew'' (1911) reports early writers who
count eight or fifteen years after the Ascension; a dispersal of the Apostles
by a tradition it calls ``admittedly not too reliable''; the Catholic critics
of its own day; a later apostolic departure after Eusebius; and St Irenaeus,
whose statement it finds too uncertain for any positive conclusion. The article ``The New
Testament'' (1912) dates the Aramaic original alone and says
nothing definite of its rendering into Greek. The Matthew article gives the first
hearers and the place: the ecclesiastical writers ``agree in declaring that St.
Matthew wrote his Gospel for the Jews'', and the Evangelist, having preached to
the Hebrews, wrote in his mother tongue ``before his departure from
Jerusalem''. The same article reports the Protestant and liberal critics of its
day as widely divided, from shortly before the fall of Jerusalem to well into
the second century. The town in the Location column locates the healing the
Gospel narrates, not the writing of the Gospel.}
\midrule
Epistle & 1 Cor 1:4--8 & Paul at Ephesus (16:8), to ``the church of God that
is at Corinth'' (1:2) & \chronodate{epistle}{\chronologyannotation{epistle}}\\
\cmidrule(lr){1-4}
\dossierprose{Traditional attribution: St Paul, with Sosthenes named beside him
(1:1). The \work{Catholic Encyclopedia}'s article on the Epistles to the
Corinthians (1908) has Paul, at Ephesus, learn of party strife and other
disorders at Corinth and of a letter from the Corinthians, write at once, and
send the Epistle by Titus; the chronological table in the article ``St.\ Paul''
(1911) gives the year of the letters to Corinth and Galatia. The two
dates are two Catholic authorities disagreeing, and both are shown. The first
article calls the historical and internal evidence for Paul's authorship so
strong that the most advanced critical schools of its day admitted it. The
same article describes Corinth as a centre of traffic and wealth, in which Paul
had spent upwards of eighteen months.}
\end{dossiertable}
